\documentclass[10pt,a4paper]{amsart}
\usepackage[margin=19mm]{geometry}
\usepackage{amsmath,amssymb,mathtools}
\usepackage[scr=rsfs,cal=euler]{mathalfa}
\usepackage{xfrac}
\usepackage[unicode,hidelinks,bookmarks=false]{hyperref}
\newcommand{\ie}{i.e.\ }
\newcommand{\cf}{cf.\ }
\newcommand{\resp}{respectively\ }
\newcommand{\Subsection}{\subsection}
\providecommand{\nobreakdash}{\nobreak\hbox{-}}
\setlength{\emergencystretch}{2em}
\multlinegap0pt
\usepackage{mathtools}
\usepackage{enumerate}
\usepackage{dsfont}
\usepackage{xcolor}
\newtheorem{lemma}{Lemma}
\newtheorem{definition}{Definition}
\newtheorem{theorem}{Theorem}
\DeclareMathOperator{\Aut}{Aut}
\newcommand{\Bir}{{\rm Bir}}
\DeclareMathOperator{\Frac}{Frac}
\newcommand{\Gal}{{\rm Gal}}
\newcommand{\PP}{\mathbb{P}}
\DeclareMathOperator{\Spec}{Spec}
\DeclareMathOperator{\id}{id}
\newcommand\rI{\mathrm{I}}
\title{Unpolarized Shafarevich conjectures for hyper-K{\"a}hler varieties}
\author{Lie Fu, Zhiyuan Li, Teppei Takamatsu, Haitao Zou}
\date{Source: arXiv:2203.10391v2 (CC BY 4.0)}
\begin{document}
\maketitle

\noindent\emph{Source and licence.} Unpolarized Shafarevich conjectures for hyper-K{\"a}hler varieties. Version arXiv:2203.10391v2 (CC BY 4.0), distributed by arXiv under the Creative Commons Attribution 4.0 International licence, http://creativecommons.org/licenses/by/4.0/, as recorded in the arXiv licence metadata for this exact version and shown on the abstract page https://arxiv.org/abs/2203.10391v2. Full licence notice: \texttt{licenses/arxiv-2203.10391-CC-BY-4.0.txt}.\par\smallskip

\noindent\textit{Editorial scope.} Counted unit: the article's lemma Lemma:Unramifiedness (source0.tex lines 1262--1287). The article's complete original proof of it is delivered with it (source0.tex lines 1289--1362, one \texttt{proof} environment of 74 lines). No mathematics is written, repaired or re-derived: the statement and the proof are the article's own lines, in the article's own notation, and no retained line is substituted or re-flowed.  Retained uncounted as the source-local context the proof needs: lines 1233--1239; lines 1903--1907; lines 1911--1927; lines 1934--1961.  Nothing was omitted from any retained span, so the package carries no omission record.  No retained line contains a citation, so the wrapper carries no bibliography and no bibliography entry.  No retained line contains a figure environment, an image-inclusion command or drawing code, so no image is delivered and the package declares no asset dependency.  Editorial formatting only, in addition: an AMS \texttt{amsart} preamble that restates the article's own declarations and macros used by the retained lines, namely the environments \texttt{lemma}, \texttt{definition}, \texttt{theorem} on separate counters, together with the standard packages those lines need.  The retained environments print in this document as Lemma 1, Definition 1, Theorem 1 in order of appearance, whereas the article prints the same items with the numbering of its own full document.  The complete native source of the article as distributed in the arXiv e-print for this exact version is bundled unchanged as \texttt{sources/123440/source0.tex}.
\begin{lemma}
\label{ruled}
Let $R$ be a discrete valuation ring with fraction field $K$ of characteristic zero.
Let $s$ be the closed point of $\Spec R$ and $k(s)$ the residue field.
Let $\mathcal{X} \rightarrow \Spec R$ be a smooth proper morphism of algebraic spaces such that the generic fiber $\mathcal{X}_{\eta}$ is a hyper-Kähler variety over $K$.
Then the special fiber $\mathcal{X}_{s}$ is a non-ruled algebraic space over $k(s)$ (Definition \ref{def:Ruled}).
\end{lemma}

\begin{definition}
\label{def:Ruled}
Let $X$ be an $n$-dimensional integral separated algebraic space of finite type over a field $k$.
We say $X$ is ruled if the schematic locus $X'$ of $X$ is ruled, i.e.,$X'$ is $k$-birationally isomorphic to $\PP^{1}_{k} \times Y$ for some scheme $Y$ of dimension $n-1$ over $k$. 
\end{definition}

\begin{theorem}[Matsusaka--Mumford for algebraic spaces]
\label{ClosedSubset}
Let $\mathcal{X}, \mathcal{Y}$ be integral smooth proper algebraic spaces over $R$.
Suppose that the special fiber $\mathcal{Y}_{s}$ is not ruled.
Let $f \colon \mathcal{X}_{\eta} \overset{\simeq}{\dashrightarrow} \mathcal{Y}_{\eta}$
be a birational map between the generic fibers.
Then there exist closed algebraic subspaces
$V \subset \mathcal{X}$, $W \subset \mathcal{Y}$
with special fibers $V_{s}, W_{s}$ satisfying $V_{s} \neq \mathcal{X}_{s}$, $W_{s} \neq \mathcal{Y}_{s}$,
such that
$f$ extends to an isomorphism over $R$
\[
  \widetilde{f} \colon 
  \mathcal{X} \setminus V \overset{\simeq}{\to}
  \mathcal{Y} \setminus W.
\]
\end{theorem}

\begin{lemma}
\label{Lemma:Specialization}
Let $\mathcal{X}$
be an integral proper smooth algebraic space over $R$.
Suppose that the special fiber $\mathcal{X}_{s}$ is not ruled.
Let $\Bir(\mathcal{X}_{\eta})$ be the group of $F$-birational automorphisms
of the generic fiber $\mathcal{X}_{\eta}$.
Let $G$ be a finite subgroup of $\Bir(\mathcal{X}_{\eta})$.
\begin{enumerate}
\item There exists a proper closed subspaces $V \subset \mathcal{X}$
with $V_s \neq \mathcal{X}_{s}$ and a homomorphism
\[
\psi \colon G \to \Aut_R(\mathcal{X} \setminus V)
\]
such that $\psi(f)|_{\mathcal{X}_{\eta}} = f$ for any $f\in G$.

\item Let $\overline{\psi}$ be the specialization morphism defined as the composition 
\[
G \overset{\psi}{\to} \Aut_R(\mathcal{X} \setminus V)
 \to \Aut((\mathcal{X} \setminus V)_{s})
\]
given by $\overline{\psi}(f) \coloneqq \psi(f)|_{(\mathcal{X} \setminus V)_{s}}$.
If the characteristic of the residue field of $R$ is $0$,
the map $\overline{\psi}$ is injective.
If the characteristic of the residue field of $R$ is $p > 0$,
the kernel of $\overline{\psi}$ is a $p$-group.
\end{enumerate}
\end{lemma}

\begin{lemma}
\label{Lemma:Unramifiedness}
Let $R$ be a henselian discrete valuation ring with $K = \Frac R$.
Let $(X, \mathscr{L})$, $(Y,\mathscr{M})$ be polarized hyper-Kähler varieties over $K$. Let
$f \colon (X_{\overline{K}}, \mathscr{L}_{\overline{K}})
\overset{\simeq}{\to}
(Y_{\overline{K}}, \mathscr{M}_{\overline{K}})$
be an isomorphism of polarized varieties over $\overline{K}$
with the associated 1-cocycle
\[ \alpha_{f} \colon \Gal(\overline{K}/K) \to \mathscr{G} \coloneqq \Aut(X_{\overline{K}}, \mathscr{L}_{\overline{K}}) \]
defined by %
$\alpha_{f} (\sigma) = f^{-1} \circ ^{\sigma} \!f$.
Assume that
\begin{itemize}
\item the order of $\mathscr{G}$ is invertible in the residue field of $R$,
\item there exist proper smooth algebraic spaces
$\mathcal{X}', \mathcal{Y}'$ over $R$
whose generic fibers are hyper-Kähler varieties
such that there are birational maps
$X \dasharrow \mathcal{X}'_{\eta}$,
$Y \dasharrow \mathcal{Y}'_{\eta}$ over $K$.
\end{itemize}
Then the restriction of $\alpha_{f}$ to the inertia subgroup
$\mathrm{I}_K \subset \Gal(\overline{K}/K)$
is trivial.
\end{lemma}

\begin{proof}
Take a finite Galois extension $L/K$ such that
$\mathscr{G} = \Aut(X_{L}, \mathscr{L}_{L})$ and 
$f$ come from an isomorphism
$f_L \colon (X_L, \mathscr{L}_{L})
\overset{\simeq}{\to}
(Y_L, \mathscr{M}_{L})$ over $L$.
The 1-cocycle  $\alpha_f$ comes from a $1$-cocycle
$$\alpha_{f,L} \colon \Gal(L/K) \to \mathscr{G}$$
corresponding to $f_L$.
Let $\rI_{L/K} \subset \Gal(L/K)$ be the inertia group.
For an element $\sigma \in \rI_{L/K}$, we have
$\alpha_{f,L}(\sigma) \coloneqq (f_L)^{-1} \circ ^{\sigma} \!f_L$ 
and 
we shall show that $\alpha_{f,L}(\sigma) = \id$.

Let $\mathcal{X}'_{\mathscr{O}_L,s}$, $\mathcal{Y}'_{\mathscr{O}_L,s}$
be the special fibers of
$\mathcal{X}'_{\mathscr{O}_L} \coloneqq \mathcal{X}' \otimes_{\mathscr{O}_K} \mathscr{O}_L$,
$\mathcal{Y}'_{\mathscr{O}_L} \coloneqq \mathcal{Y}' \otimes_{\mathscr{O}_K} \mathscr{O}_L$,
respectively.
We note that the isomorphism $f_L \colon X_L \overset{\simeq}{\to} Y_L$
can be considered as a birational map
\[f_L \colon \mathcal{X}'_{L} \overset{\simeq}{\dashrightarrow} \mathcal{Y}'_{L}.\]
Here, we put $\mathcal{X}'_{L} \coloneqq \mathcal{X}' \otimes_{\mathscr{O}_K} L$ and
$\mathcal{Y}'_{L} \coloneqq \mathcal{Y}' \otimes_{\mathscr{O}_K} L$. By Theorem \ref{ClosedSubset} and Lemma \ref{ruled},
the birational map $f_L$
extends to an isomorphism
\[
  \widetilde{f}_L \colon
  (\mathcal{X}'_{\mathscr{O}_L} \setminus V) \overset{\simeq}{\to}
  (\mathcal{Y}'_{\mathscr{O}_L} \setminus W)
\]
over $\mathscr{O}_L$
for some proper closed subsets
$V \subset \mathcal{X}'_{\mathscr{O}_L}$, $W \subset \mathcal{Y}'_{\mathscr{O}_L}$
satisfying
$V_s \neq \mathcal{X}'_{\mathscr{O}_L,s}$, $W_s \neq \mathcal{Y}'_{\mathscr{O}_L,s}$.
Enlarging $V$ and $W$ if necessary, by Lemma \ref{Lemma:Specialization},
we may assume the following holds:
The action of
$\mathscr{G} = \Aut(X_{L}, \mathscr{L}_{L})$ on $X_L$
extends to a homomorphism
\[
  \psi \colon \mathscr{G} \to \Aut(\mathcal{X}'_{\mathscr{O}_L} \setminus V).
\]

Since $\sigma \in \rI_{L/K}$ acts trivially on the residue field of $\mathscr{O}_L$,
the isomorphisms
\[
  \widetilde{f}_L \colon
  (\mathcal{X}'_{\mathscr{O}_L} \setminus V) \overset{\simeq}{\to}
  (\mathcal{Y}'_{\mathscr{O}_L} \setminus W)
  \quad \text{and} \quad
  ^{\sigma} \!\widetilde{f}_L \colon
  (\mathcal{X}'_{\mathscr{O}_L} \setminus \sigma(V)) \overset{\simeq}{\to}
  (\mathcal{Y}'_{\mathscr{O}_L} \setminus \sigma(W))
\]
induce the same morphism on the special fiber.
Thus, the element
\[
 \alpha_{f,L}(\sigma) = (f_L)^{-1} \circ ^{\sigma}\!f_L \in \mathscr{G} \subset \Bir (\mathcal{X}_{L})
 \]
sits in the kernel of the specialization map
\[
  \overline{\psi} \colon \mathscr{G} \to \Aut((\mathcal{X}'_{\mathscr{O}_L} \setminus V)_s).
\]
Since the order of $\mathscr{G}$ is invertible
in the residue field of $R$,
the specialization map
$$\mathscr{G} \to \Aut((\mathcal{X}'_{\mathscr{O}_{L}} \setminus V)_{s})$$
is injective by Lemma \ref{Lemma:Specialization} (2).
Therefore,  we have $\alpha_{f,L}(\sigma) = \id$.
\end{proof}

\end{document}
